% GraphMend, CGO 2027 -- Artifact Appendix.
%
% Structure follows the ctuning AE template (ae-20190108.tex); style and
% length follow the PyTorch 2 artifact appendix (ASPLOS '24), which is the
% closest comparable in this space.
%
% Place BEFORE \bibliography. Requires \usepackage{artifact-evaluation} or
% equivalent; plain \appendix works too.
%
% PLACEHOLDERS TO FILL: {{REPO_URL}}, {{COMMIT}}, {{DOI}}.

\appendix
\section{Artifact Appendix}

\subsection{Abstract}

The artifact contains the GraphMend implementation, its three AST-level
transformations (Predicated Data-Dependent Control Flow, Graph-Epilogue
Deferred Side Effects, and Predicated Trap Lowering), the per-model benchmark
drivers for the 27 Hugging Face models of Table~\ref{tab:benchmark}, and
scripts for graph-break detection, output-correctness checking, full-graph
capture, latency and throughput measurement, profiler collection, and
compiler-overhead measurement. It reproduces the key results of
Section~\ref{sec:eval}.

The primary claims -- break elimination, output correctness, and full-graph
capture -- are \emph{CPU-only}: they need no GPU, no pretrained weights, and,
for 21 of the 27 models, no network access. The latency and throughput results
of Sections~\ref{sec:coldstart}--\ref{sec:throughput} depend on the GPUs used
in the paper and are provided as scripts that measure on the reviewer's own
hardware. A container image pins the software stack, and a single
\texttt{run\_quick.sh} gives a pass/fail functional result in 10--20 minutes.

\subsection{Artifact check-list (meta-information)}

{\small
\begin{itemize}
  \item {\bf Algorithm:} three source-level, legality-guarded program
    transformations applied to PyTorch model code claimed at import time.
  \item {\bf Program:} GraphMend, implemented in the Jaseci/Jac compiler;
    PyTorch 2.12.1; Hugging Face Transformers 4.52.4.
  \item {\bf Compilation:} the toolchain is upstream \texttt{jaseci-labs/jaseci}
    at commit \texttt{e2b6b9f4b} (jaclang 0.36.1) plus
    \texttt{patches/graphmend.patch} (203 files, 166 new). It is used from
    source on \texttt{PYTHONPATH} and declares no runtime PyPI dependencies.
  \item {\bf Transformations:} Predicated Data-Dependent Control Flow
    (\textsc{[Where]}), Graph-Epilogue Deferred Side Effects (\textsc{[Defer]}),
    Predicated Trap Lowering (\textsc{[Trap]}); enabled by
    \texttt{[run] graphmend = true} (default) together with
    \texttt{[run] graphmend\_claim\_imports = true} (\emph{not} default).
  \item {\bf Model:} the 27 Hugging Face models of Table~\ref{tab:benchmark}
    that exhibit graph breaks. 21 are built offline from small configurations
    with random weights; 6 require network access and
    \texttt{trust\_remote\_code}.
  \item {\bf Data set:} none. Inputs are synthetic tensors of fixed shape,
    seeded with \texttt{torch.manual\_seed(0)}.
  \item {\bf Run-time environment:} Linux or macOS, Python 3.13. Container
    images pin the full stack for both the CPU and the GPU experiments.
  \item {\bf Hardware:} any x86-64 or arm64 CPU with 12\,GB of memory for the
    primary claims. Latency and throughput need an NVIDIA GPU; the paper used
    RTX 3090, A40 and H100 with CUDA 12.6 and Triton 3.7.
  \item {\bf Run-time state:} cold-cache sensitive. Each arm of the latency
    benchmark is given a private TorchInductor cache so that cold start is
    genuinely cold; the container warms the Jac compiler cache at build time.
  \item {\bf Execution:} \texttt{bash scripts/run\_quick.sh} for the functional
    workflow; \texttt{bash scripts/run\_full.sh} for the full reproduction.
  \item {\bf Metrics:} graph-break count (FX graphs $-$ 1); fix rate; output
    fingerprint equality; whether \texttt{torch.compile(fullgraph=True)}
    captures; cold-start and steady-state forward-pass latency; CUDA-graph
    launches per forward; end-to-end throughput; compiler overhead over
    standard Python.
  \item {\bf Output:} a table per experiment, printed to stdout, with the
    claim and its expected result stated in the banner of each script.
    Reference output is in \texttt{artifact/README.md}.
  \item {\bf Experiments:} 4 rule-level graph-count suites (18 tests); 21
    offline model rows; 6 network model rows; per-break cause reporting;
    full-graph capture per arm; compiler overhead; and, on a GPU, break counts,
    CUDA-graph launches, cold start, steady state and throughput.
  \item {\bf How much disk space required (approximately)?:} 135\,MB for a
    clone with the submodule; 1.4\,GB for the built CPU image, or roughly
    5\,GB with the build cache. The 6 network rows add a few GB of Hub
    downloads.
  \item {\bf How much time is needed to prepare workflow (approximately)?:}
    about 5 minutes for the container build, most of it the PyTorch wheel
    download and the build-time compiler-cache warm.
  \item {\bf How much time is needed to complete experiments
    (approximately)?:} 10--20 minutes for the functional workflow; roughly
    1.5--3 hours for the 21-row offline sweep on a cold cache; a further
    30--60 minutes for the GPU experiments.
  \item {\bf Publicly available?:} yes, {{REPO\_URL}} at commit {{COMMIT}}.
  \item {\bf Code licenses (if publicly available)?:} MIT.
  \item {\bf Data licenses (if publicly available)?:} not applicable; no data
    set is distributed.
  \item {\bf Workflow framework used?:} none; plain shell and Python.
  \item {\bf Archived (provide DOI)?:} {{DOI}}.
\end{itemize}
}

\subsection{Description}

\subsubsection{How to access}

\begin{itemize}
  \item Repository: {{REPO\_URL}}, commit {{COMMIT}}. Clone with
    \texttt{-{}-recurse-submodules}.
  \item Archived copy: {{DOI}}. The deposit is self-contained: the toolchain
    and the typeshed stubs are flattened into it, so it needs no network.
\end{itemize}

The compiler is not vendored. \texttt{jaseci/} is the upstream toolchain as a
submodule frozen at \texttt{e2b6b9f4b}, and
\texttt{patches/graphmend.patch} is everything GraphMend adds to it, so the
patch is a direct statement of the contribution.
\texttt{scripts/setup.sh} verifies the pin, applies the patch, and fetches the
typeshed stubs; it is idempotent.

\subsubsection{Hardware dependencies}

\begin{itemize}
  \item Break elimination, correctness, full-graph capture and compiler
    overhead: any x86-64 or arm64 CPU. \textbf{12\,GB of memory is required};
    \texttt{Phi-4-mini-instruct} peaks near 9.7\,GB while GraphMend compiles
    the imported modeling code, and a smaller ceiling causes the container to
    be killed by the OOM killer.
  \item Latency and throughput: one NVIDIA GPU. Functional validation works on
    any sufficiently recent device; reproducing the magnitudes of
    Table~\ref{tab:benchmark} requires the RTX 3090, A40 or H100
    configurations used in the paper.
\end{itemize}

\subsubsection{Software dependencies}

\begin{itemize}
  \item A recent Linux distribution or macOS, and \texttt{git}
  \item Python 3.13
  \item PyTorch 2.12.1, Transformers 4.52.4, NumPy 2.4.6, torchvision 0.27.1
  \item For the GPU experiments: NVIDIA drivers, CUDA 12.6, Triton 3.7
  \item Docker, for the pinned container path
\end{itemize}

Versions are pinned exactly rather than floored. A graph-break count is a
property of what TorchDynamo chooses to split on, so a different PyTorch
release can report a different number, which a reviewer could reasonably read
as a failed reproduction.

\subsubsection{Data sets}

None. Inputs are synthetic tensors at fixed shapes, seeded so that the two arms
receive identical weights and identical inputs. The harness prints the input
shape per arm and flags a mismatch rather than assuming one.

\subsubsection{Models}

The 27 rows of Table~\ref{tab:benchmark}. 21 are built offline from small
configurations with random weights; graph breaks are structural, so this
preserves the measured quantity. 6 need network access and
\texttt{trust\_remote\_code}: \texttt{Florence-2},
\texttt{MoLFormer-XL-both10pct} (revision pinned),
\texttt{chronos-bolt-small}, \texttt{Qwen-Audio-Chat},
\texttt{stella-en-400M-v5} and \texttt{moe-minicpm-x4-base}.

\subsection{Installation}

\begin{verbatim}
git clone --recurse-submodules {{REPO_URL}}
cd CGO_AE_GraphMend
docker build -f artifact/Dockerfile.cpu \
    -t graphmend-cpu .
docker run --rm graphmend-cpu
\end{verbatim}

Or natively, with the dependencies above installed:

\begin{verbatim}
bash scripts/setup.sh
bash scripts/run_quick.sh
\end{verbatim}

\subsection{Experiment workflow}

Two workflows are provided. The functional workflow answers whether GraphMend
works, in 10--20 minutes on a CPU with no network:

\begin{verbatim}
bash scripts/run_quick.sh
\end{verbatim}

It runs the four rule-level graph-count suites, which exercise each rule on
fixtures, and then four representative models: \texttt{t5-small}
(\textsc{[Defer]}), \texttt{Phi-4-mini-instruct} (\textsc{[Where]}, the
Figure~\ref{fig:phi4} worked example), \texttt{longformer-base-4096} and
\texttt{clap-htsat-fused}. The last two are expected \emph{not} to reach zero
breaks; a run that fixed them would indicate a problem rather than a success.
It prints \texttt{PASS}/\texttt{FAIL} per check and exits non-zero on failure.

The full reproduction runs every claim the available hardware allows, skipping
the two GPU experiments where no CUDA device is visible:

\begin{verbatim}
bash scripts/run_full.sh
\end{verbatim}

Two properties of the setup are easy to get wrong, and both produce a clean run
that silently measures nothing. First, \texttt{graphmend\_claim\_imports} is off
by default and all model code is imported third-party code, so without the
opt-in both arms are identical. Second, the entry program must be compiled by
Jac: GraphMend injects the deferred-side-effect hooks at the
\texttt{torch.compile(...)} assignment site, so under plain CPython no hook is
registered and every logger break survives. Both are documented prominently in
the artifact README, and both are handled by the shipped scripts.

\subsection{Evaluation and expected results}

Each claim has one script.

\begin{itemize}
  \item \texttt{run\_break\_analysis.sh} $\rightarrow$
    Table~\ref{tab:benchmark}. Expect \textbf{21 models fully fixed, 3
    partially fixed and 3 unfixed}, matching the table's model-level outcome,
    with 25 of 27 rows matching its fix rate. \texttt{grounding-dino} reports
    56\% against 58\%, a difference of one break out of seventeen.
  \item \texttt{run\_correctness.sh} $\rightarrow$
    Section~\ref{sec:eval}. Expect \texttt{output\_ok} to read \texttt{yes} on
    every row, including the rows that fix nothing: FP32 outputs are
    bit-identical and greedy-decoded token sequences agree.
  \item \texttt{run\_fullgraph.sh} $\rightarrow$ Section~\ref{sec:serving}.
    Expect \texttt{torch.compile(fullgraph=True)} to \emph{fail} on every
    original model and \emph{succeed} on every transformed one. The check is
    asymmetric by construction: both-pass or both-fail is a failure.
  \item \texttt{run\_latency.sh} $\rightarrow$ Table~\ref{tab:benchmark},
    Sections~\ref{sec:coldstart}--\ref{sec:steadystate}. Expect graph breaks to
    reach zero and CUDA-graph launches per forward to collapse to one, and a
    cold-start speedup well above unity.
  \item \texttt{run\_throughput.sh} $\rightarrow$ Figure~\ref{fig:throughput}.
    Expect a gain that tracks the number of CUDA-graph launches eliminated and
    that shrinks as the batch grows.
  \item \texttt{run\_compiler\_overhead.sh} $\rightarrow$
    Figure~\ref{fig:overhead}, Section~\ref{sec:overhead}. Expect a mean near
    11.5\% on a cold compiler cache and near 1.1\% on a warm one.
\end{itemize}

\paragraph{Expected variation.} The deterministic results must match exactly:
graph-break counts, CUDA-graph launch counts, output fingerprints, and
full-graph capture outcomes. Timing results vary with the machine. Cold start
is validated against a floor of $1.5\times$ on any CUDA device rather than
against an expected value; steady state and throughput are reported without a
validation threshold, because both move with the card and the batch size and a
fixed bound would fail an honest run on different hardware.

\paragraph{On absolute break counts.} The CPU sweep runs in FP32, which is what
the correctness claim concerns, and its absolute counts therefore differ from
those in Table~\ref{tab:benchmark}, which were measured in half precision on a
GPU: 122 breaks against 147, with 17 of 26 rows matching exactly. The largest
group that differs is the BART family, whose overflow guard leads with
\texttt{hidden\_states.dtype == torch.float16}. That conjunct is a static
Python bool, so in FP32 Dynamo folds it to \texttt{False} and the
data-dependent breaks behind it do not exist. Measured in half precision
through \texttt{artifact/gpu/bench.py -{}-count}, those rows match the table
exactly. The fix rate, which is what the \emph{Fixed(\%)} column reports, is
unaffected.

\subsection{Experiment customization}

Individual rows run by name, for example
\texttt{bash scripts/run\_break\_analysis.sh t5-small biogpt}, and
\texttt{GM\_GPU\_MODELS} selects a subset for the GPU scripts. Individual rules
can be ablated with \texttt{[run] graphmend\_disable}, which takes a
comma-separated list of \texttt{trap}, \texttt{where} and \texttt{defer};
setting \texttt{graphmend = false} disables all three.

To point GraphMend at a model that is not in the suite, add a \texttt{\_build()}
returning \texttt{(model, inputs)} and a registry entry naming the modules
GraphMend should transform, then run the analyzer to see the detected breaks
and their causes:

\begin{verbatim}
python -m paper_eval.run_why <model_key> on
\end{verbatim}

This prints each surviving break with its reason and source location, which is
what distinguishes a correctly unfixed row from a defect. Note that a model's
breaks frequently live in shared \texttt{transformers} modules rather than in
its own package, so a transform scope naming only the model package will
silently miss them.

\subsection{Notes}

Enabling the code path under test matters as much as enabling GraphMend.
\texttt{Phi-4-mini-instruct} only takes the LongRoPE branch when
\texttt{rope\_scaling.type} is \texttt{"longrope"}; \texttt{clap} needs
\texttt{enable\_fusion}; \texttt{longformer} needs a sequence length that is not
a multiple of the attention window; and \texttt{Florence-2}'s guard is dead code
in FP32 and must be measured in half precision. With the default configuration
each of those rows would pass vacuously.

\texttt{stella-en-400M-v5} reproduces its 0\% row only on CUDA with xformers
installed, because its breaks live behind the unpadding path. The CPU fallback
measures a different, smaller break set.

\subsection{Methodology}

Submission, reviewing and badging methodology:

\begin{itemize}
  \item \url{https://www.acm.org/publications/policies/artifact-review-badging}
  \item \url{http://cTuning.org/ae/submission-20201122.html}
  \item \url{http://cTuning.org/ae/reviewing-20201122.html}
\end{itemize}
